% Luis
%=============================================
% Puntos 1-4: Resultado final, cercanía fuente estimada/real,
% relación distancia-amplitud, coherencia del modelo.
%=============================================

\section{Conclusiones y Defensa Final}

\begin{frame}{Resultados Finales: Validación del Modelo}

    \begin{columns}[T]

        %===========================================
        % COLUMNA IZQUIERDA: Comparación numérica
        %===========================================
        \column{0.50\textwidth}

        \begin{block}{Fuente Real vs. Estimada}
            \centering
            \scriptsize
            \begin{tabular}{@{}lccc@{}}
                \toprule
                & $x$ & $y$ & $z$ \\ \midrule
                \textbf{Real}      & $1.200$ & $-0.800$ & $-1.100$ \\[0.15cm]
                \textbf{Estimada}  & $1.200$ & $-0.777$ & $-1.138$ \\[0.15cm] \midrule
                \textbf{Error} $|\Delta|$ & $0.000$ & $0.023$ & $0.038$ \\
                \bottomrule
            \end{tabular}
            \vspace{0.15cm}
            \tiny \textcolor{gray}{Coordenadas en kilómetros (km).}
        \end{block}

        \vspace{0.2cm}

        \begin{exampleblock}{Resultado del Proyecto}
            \scriptsize
            El proyecto estimó correctamente la ubicación de la fuente sísmica a partir de las amplitudes registradas por sensores distribuidos en el espacio.
        \end{exampleblock}

        %===========================================
        % COLUMNA DERECHA: Interpretación
        %===========================================
        \column{0.46\textwidth}

        \begin{block}{Cercanía e Interpretación Física}
            \scriptsize
            \begin{itemize}
                \item La fuente estimada quedó \textbf{muy cerca} de la fuente real, validando el modelo.
                \vspace{0.1cm}
                \item La amplitud \textbf{disminuye con la distancia} por atenuación y dispersión de la energía.
                \vspace{0.1cm}
                \item Las regiones de menor error \textbf{coinciden} con posiciones cercanas a la fuente real.
            \end{itemize}
        \end{block}

        \vspace{0.2cm}

        \begin{block}{Coherencia del Modelo}
            \centering
            \scriptsize
            El comportamiento físico y matemático obtenido es \textbf{consistente} con lo esperado.
        \end{block}

    \end{columns}

    \vspace{0.15cm}
    \begin{center}
        \footnotesize
        \textbf{El modelo reprodujo adecuadamente las mediciones observadas.}
    \end{center}

\end{frame}